

\section{Linear algebra warm-up: PC for matrices}
\label{sec:matrix}
Two matrices $M,M' \in U \otimes V$ are isomorphic as tensors if they are equivalent as matrices, meaning under left- and right-multiplication by invertible matrices $X \in \GL(U), Y \in GL(V)$, that is, 
\[
XMY = M'. 
\]

Since we want $X,Y$ to be invertible, we also introduce variable matrices $X', Y'$ as before, together with the equations
\[
XX' = X'X = \Id_U \qquad YY' = Y'Y = \Id_V.
\]
Then by left multiplying our initial matrix equation by $Y'$, we may replace it with the new matrix equation
\[
XM = M'Y'.
\]
The latter has the advantage of being linear in $X$ and $Y'$, but the quadratic equations $XX' =  \Id_U, YY' =  \Id_V$ still make even this case not totally obvious.

\subsection{A trick for PC degree} \label{sec:trick}
If our focus is on PC \emph{degree}, we note that the degree of the equations is unchanged if we first left- or right-multiply $M,M'$ by invertible scalar matrices. For example, if we replace $M$ by $\overline{M} = AMB$ with $A,B \in \GL(U)$, then we may replace $X$ by $\overline{X} := XA^{-1}$, $Y$ by $\overline{Y} := B^{-1} Y$. Then we have $\overline{M} \cong M$, so $\overline{M} \cong M'$ iff $M \cong M'$. Furthermore, since the transformation $X \mapsto XA^{-1}$, $Y \mapsto B^{-1} Y$ is linear and invertible, any PC proof that $\overline{M} \not\cong M'$ can be transformed by the inverse linear transformation into a PC proof that $M \not\cong M'$ of the same degree.

Now, for matrices under this equivalence relation, we have a normal form, namely every matrix $M$ is equivalent to a diagonal matrix with $\rk(M)$ 1s on the diagonal and all the remaining entries $0$, that is, $\sum_{i=1}^{\rk(M)} e_i \otimes e_i = I_r \oplus 0$, where the latter $0$ denotes a $0$ matrix of appropriate size $(n-r) \times (m-r)$. So by using the preceding trick, we may put both $M$ and $M'$ in this form. The two are isomorphic iff $\rk(M) = \rk(M')$, so for PC degree we have now reduced to the case of showing that $I_r \oplus 0$ and $I_{r'} \oplus 0$ are not isomorphic when $r \neq r'$. 

Note that, aside from the equations saying $X$ and $Y$ are invertible, this is almost identical to the Rank Principle (see Section~\ref{sec:prelim:inv}). In the rest of this section we will prove PC lower bounds on both the Rank Principle and the Inversion Principle. Here, we show that the addition of these invertibility axioms in fact makes 2TI much easier in PC than the Rank Principle or 3TI.

\begin{proposition}
Let $M,M'$ be two $n \times m$ matrices of  ranks $r,r'$ respectively, with $r' > r$. Then, over any field whose characteristic does not divide $r'-r$, the following equations have a degree 3 PC refutation and a degree 4 NS refutation:
\[
XMY^T = M' \qquad XX' = X' X = \Id_n \qquad YY' = Y' Y = \Id_m.
\]
\end{proposition}

For those familiar with the low-degree PC proof of the functional onto-PHP, the following proof is similar. 
\newcommand{\Tr}{\text{Tr}}

\begin{proof}[Proof idea]
By the observations in Section~\ref{sec:trick}, we may assume without loss of generality (from the point of view of PC degree) that $M = \Id_r \oplus 0_{n - r \times m-r}$ and $M' = \Id_{r'} \oplus 0_{n-r' \times m-r'}$. 

Write $X = \begin{bmatrix} X_{11} & X_{12} \\ X_{21} & X_{22} \end{bmatrix}$ where the top-left block $X_{11}$ has size $r' \times r$, and similarly write $Y = \begin{bmatrix} Y_{11} & Y_{12} \\ Y_{21} & Y_{22} \end{bmatrix}$ where $Y_{11}$ has size $r' \times r$. In this notation, the matrix equation $XMY^T = M'$ becomes the equations
\begin{align*}
XMY^T & = \begin{bmatrix} X_{11} & X_{12} \\ X_{21} & X_{22} \end{bmatrix} \begin{bmatrix} \Id_r \\ & 0_{(n-r) \times (m-r)} \end{bmatrix} \begin{bmatrix} Y_{11}^T & Y_{21}^T \\ Y_{12}^T & Y_{22}^T  \end{bmatrix} \\
 & = \begin{bmatrix} X_{11} & 0 \\ X_{21} & 0 \end{bmatrix} \begin{bmatrix} Y_{11}^T & Y_{21}^T \\ 0 & 0 \end{bmatrix} \\
 & = M' = \begin{bmatrix} \Id_{r'} \\ & 0_{(n-r') \times (m-r')} \end{bmatrix}
\end{align*}
which becomes the four matrix equations
\begin{equation} \label{eq:2}
X_{11} Y_{11}^T = \Id_{r'} \qquad X_{11}Y_{21}^T = 0 \qquad X_{21}Y_{11}^T = 0 \qquad X_{21}Y_{21}^T = 0.
\end{equation}
Note that so far our PC proof hasn't actually done anything---it is all just notation, and all in the same degree we started with (degree 2).

Then, using the equations $XX' = \Id_n$ and $YY'=\Id_m$, we will derive that $Y_{11}^T X_{11} = \Id_r$. Then we derive 1 as 
\[
\frac{1}{r-r'}\left(\Tr(X_{11} Y_{11}^T - \Id_{r'}) - \Tr(Y_{11}^T X_{11} - \Id_r)\right).
\]
The point here is that trace is additive and cyclically invariant, so $\Tr(X_{11} Y_{11}^T) \equiv \Tr(Y_{11}^T X_{11})$, identically as polynomials, so there is no further derivation needed.
\end{proof}

\begin{proof}
The proof starts using the first part of the proof idea above, so we continue from Equation~(\ref{eq:2}) with the notation introduced above.
In the remainder of the proof, we will derive $Y_{11}^T X_{11} = \Id_r$. Then the last paragraph of the proof idea will complete the proof.

To derive $Y_{11}^T X_{11} = \Id_r$, we will use the invertibility equations (those involving $X'$ and $Y'$). Write $X' = \begin{bmatrix} X'_{11} & X'_{12} \\ X'_{21} & X'_{22} \end{bmatrix}$, where $X'_{11}$ has size $r \times r'$ (NB: the size is the ``transpose'' of the size of $X_{11}$) and similarly for $Y'$. 

From considering the upper-left $r \times r$ block of the matrix equation $X' X = \Id_n$, we get
\[
X'_{11} X_{11} + X'_{12} X_{21} = \Id_r.
\]
Right multiplying by $Y_{11}^T$, we get
\[
X'_{11} X_{11} Y_{11}^T + X'_{12} X_{21} Y_{11}^T = Y_{11}^T. 
\]
But now we can subtract from this $X'_{11}$ times the equation $X_{11} Y_{11}^T = \Id_{r'}$, and also $X'_{12}$ times the equation $X_{21} Y_{11}^T = 0$ to get
\begin{equation}
X'_{11} = Y_{11}^T.
\end{equation}

Similarly, considering the upper-left $r \times r$ block of the matrix equation $Y' Y = \Id_m$, we get $Y'_{11} Y_{11} + Y'_{12} Y_{21} = \Id_{r}$. For consistency with the notation above, we take the transpose of this entire equation (in PC, this is essentially a null-op---we are just re-arranging how we are viewing a set of $(r')^2$ equations on the page), to get:
\[
Y_{11}^T (Y'_{11})^T + Y_{21}^T (Y'_{12})^T = \Id_r.
\]
Left multiplying by $X_{11}$, we get
\[
X_{11} Y_{11}^T (Y'_{11})^T + X_{11} Y_{21}^T (Y'_{12})^T = X_{11}.
\]
Now, right-multiplying the equation $X_{11} Y_{11}^T = \Id$ by $(Y'_{11})^T$, and right-multiplying the equation $X_{11} Y_{21}^T = 0$ by $(Y'_{12})^T$ and subtracting both of these from the above, we get
\begin{equation}
(Y'_{11})^T = X_{11}.
\end{equation}

Next, we derive $M - X' M' (Y')^T=0$, as follows: left-multiplying $XMY^T - M'$ by $X'$, and subtract from it $(X'X-I)$ times $MY^T$, to get $-X'M' + MY^T$. Now right-multiply the latter by $(Y')^T$ and subtract from it $M$ times $(Y^T (Y')^T - I)$, yielding $-X'M'(Y')^T + M$. Now multiply by $-1$.

Now, from $X' M' (Y')^T = M$, as at the beginning of the proof, we derive that $X'_{11} (Y'_{11})^T = \Id_r$. But above we have derived that $X'_{11} = Y_{11}^T$ and $(Y'_{11})^T = X_{11}$, so from the preceding three equations we get $Y_{11}^T X_{11} = \Id_r$, as claimed. This completes the PC proof.

Let us unroll the PC proof to derive a Nullstellensatz proof (here we underline the use of original equations):
\begin{align*}
r - r' & = \Tr(\underline{X_{11} Y_{11}^T - \Id_{r'}}) - \Tr(Y_{11}^T X_{11} - \Id_r)  \\
\end{align*}
Now we focus on the NS derivation of $Y_{11}^T X_{11} - \Id_r$. Since the trace is linear, and we are focusing on degree, this is without loss of generality. We have:
\begin{align*}
Y_{11}^T X_{11} - \Id_r = & \left(X'_{11} (Y'_{11})^T - \Id_r\right) - (X'_{11} - Y_{11}^T) (Y'_{11})^T -  Y_{11}^T \left((Y'_{11})^T - X_{11}\right) \\
 = & \left(-X' \underline{(XMY^T - M')} Y'^T + \underline{(X'X - I)}MY^T (Y')^T + M \underline{(Y^T (Y')^T - I)}\right)_{11} \\
  & + \left(\underline{(X'_{11} X_{11} + X'_{12} X_{21} - \Id_r)} Y_{11}^T - X'_{11} \underline{(X_{11} Y_{11}^T - \Id_{r'})} - X'_{12}\underline{(X_{21} Y_{11}^T)} \right) (Y'_{11})^T \\ 
  & -  Y_{11}^T \left( X_{11} \underline{(Y_{11}^T (Y'_{11})^T + Y_{21}^T (Y'_{12})^T - \Id_r)} - \underline{(X_{11} Y_{11}^T - \Id)} (Y'_{11})^T - \underline{(X_{11} Y_{21}^T)} (Y'_{12})^T   \right). 
\end{align*}
This is visibly degree 4.
\end{proof}

\subsection{Inversion Principle implies the Rank Principle}

\begin{lemma} \label{lem:inv_to_rank}
If the $r \times r$ Inversion Principle has a degree $d$ PC derivation, 
then there is a degree $\max\{d,3\}$ PC refutation of the Rank Principle stating that a rank $r$ matrix is not equivalent (isomorphic) to a rank $n$ matrix, for any $n > r$. 

If the Inversion Principle has a degree $d$ NS derivation, then the Rank Principle has a degree $d+2$ NS refutation.
\end{lemma}

\begin{proof}
Suppose the $r \times r$ Inversion Principle has a degree-$d$ derivation. Consider the Rank Principle $XY=I_n$ where $X$ is $n \times r$ and $Y$ is $r \times n$, with $n > r$. Write
\[
X = \begin{bmatrix} X_0 \\ X_1 \end{bmatrix} \text{ and } Y = \begin{bmatrix} Y_0 & Y_1 \end{bmatrix},
\]
where $X_0,Y_0$ are $r \times r$. Then, examining the upper-left $r \times r$ corner of the original equations, we find $X_0 Y_0 = I_r$. As these are square matrices, by assumption in degree $d$ we may then derive that $Y_0 X_0 = I_r$ as well.

Now, multiply both sides of $XY=I_n$ on the left by the matrix $\begin{bmatrix} Y_0 & 0 \\ 0 & I_{n-r} \end{bmatrix}$. The result is then the set of degree-3 equations
\[
\begin{bmatrix} Y_0 X_0 \\ X_1 \end{bmatrix} \begin{bmatrix} Y_0 & Y_1 \end{bmatrix} = \begin{bmatrix} Y_0 & 0 \\ 0 & I_{n-r} \end{bmatrix}.
\]
Considering the upper-right $r \times (n-r)$ block of these equations, we find the equations $Y_0 X_0 Y_1 = 0$. 

But now, from the equation $Y_0 X_0 = I_r$, we may right-multiply by $Y_1$ to get $Y_0 X_0 Y_1 = Y_1$. Combining with the equation at the end of the last paragraph, we then conclude $Y_1 = 0$.

Finally, consider the lower-right $(n-r) \times (n-r)$ part of the original equation $XY = I_n$, namely, $X_1 Y_1 = I_{n-r}$. We had already derived $Y_1 = 0$, which we can then left-multiply by $X_1$ to get $X_1 Y_1 = 0$. Considering any diagonal entry of these two equations, we then derive the contradiction $1=0$.

To see the NS certificate, we unwrap the above proof. First write $Y_0 X_0 - I_r$ as a linear combination of the equations $X_0 Y_0 - I_r$ with polynomial coefficients, in total degree $d$. Among our starting equations in the Rank Principle, we have $X_0 Y_1$ and $X_1 Y_1 - I_{n-r}$. Then the following linear combination has degree 2 more than $Y_0 X_0 - I_r$, and derives $1$ in any of its diagonal entries: 
\[
-X_1 Y_0 \underline{X_0 Y_1} + X_1 (Y_0 X_0 - I_r) Y_1 + \underline{(X_1 Y_1 - I_{n-r})}.
\]
\end{proof}

\begin{observation}
The $n \times n$ Inversion Principle has a proof of degree $2n+2$.
\end{observation}

\begin{proof}
The idea is to use Laplace expansion. We spell out the details.

We start with $XY = I_n$, where $X$ and $Y$ are $n \times n$ matrices of variables. Left-multiply by $Y$ to get $YXY = Y$, and then right multiply by $Adj(Y)$ (whose entries are the $(n-1) \times (n-1)$ cofactors of $Y$, hence have degree $n-1$) to get $YXYAdj(Y) = YAdj(Y)$. Now, by Laplace expansion, we have $Y Adj(Y) \equiv \det(Y) I_n$, so we get $YX \det(Y) = \det(Y) I_n$.

Next, starting from $XY=I_n$ and expanding out the determinant term-by-term, we derive $\det(XY) = 1$. (Note that here, we are not simply applying the determinant to the matrix $XY-I$, as that would give us the value of the characteristic polynomial evaluated at 1. Instead, we repeatedly use that from $a-b=0$ and $c-d=0$ we can derive $ac-bd=0$ as $\underline{(a-b)}c + b\underline{(c-d)}$. Similarly, we can derive $(a+c)-(b+d)=0$ as $(a-b) + (c-d)$.) Now, since $\det(XY) \equiv \det(X)\det(Y)$ identically as polynomials, we have derived $\det(X)\det(Y)=1$ in degree $n$.

Now, from $YX \det(Y) - \det(Y) I_n$ in the first paragraph, we multiply by $\det(X)$ to get $(YX - I_n)(\det(X)\det(Y))$. From $\det(X)\det(Y)-1$ in the second paragraph, we multiply by $-(YX-I_n)$ and add to the preceding to get $YX - I_n$, all in degree at most $2n+2$.
\end{proof}


\subsection{Lower bound on the Rank Principle (and Inversion Principle) via reduction from PHP}
\newcommand{\bbI}{\mathbb{I}}

Here we show that the Rank Principle (see Section~\ref{sec:prelim:inv}) requires large PC degree, via a reduction to the Pigeonhole Principle. 
For the Pigeonhole principle, a tight PC degree lower bound is known:


\begin{theorem}[{Razborov \cite{DBLP:journals/cc/Razborov98}}]
\label{thm:PCdeglb}
Any $PC$ refutation of the Functional $PHP^n_r$ requires degree $r/2$+1 over any field.
\end{theorem}

We use this to show:

\begin{theorem} \label{thm:rank}
Let $n\in \mathbb N$,  $n\geq 2$ and $1\leq r<n$. $\bbI(r,n)$ (with or without the Boolean axioms) requires degree $r/2+1$ in $PC$ over any field.
\end{theorem}
\begin{proof}
We prove that $PHP^n_{r}$ is $(1,2)$-reducible to $\bbI(r,n)$. 
First we consider the following degree $1$  polynomials  defining $x$ and $y$ variables of $\bbI(r,n)$ 
in terms of the $p$ variables of $PHP^n_{r}$.
variables 
$$
x_{i,k} = y_{i,k} = p_{i,k} \qquad \mbox{ for } i \in [n], k \in [n-1].
$$ 
Second we show a degree $2$ $PC$ proof   of $\bbI(r,n)$ from the  polynomials defining the  $PHP^n_r$.
From $PHP$ axioms $p_{i,k}p_{k,j}$ for $i,j \in [n], i\not =j$,  and summing over all $k \in [r]$, we get 
$$
\sum_{k \in [r]}p_{i,k}p_{k,j},
$$ 
which are exactly the axioms of $\bbI(r,n)$ for $i \not =j, i,j \in [n]$, after the substitution of variables.

For a $i\in [n]$, take  the boolean axioms written in the form $p_{i,k}p_{i,k}-p_{i,k}$ and  sum them over $k \in [r]$:
$$\sum_{k \in [r]} p_{i,k}p_{i,k} - \sum_{k \in [r]}p_{i,k}$$
Summing this last polynomial with the $PHP$ axiom $\sum_{k \in [r]} p_{i,k} -1$
we get the polynomial
$$\sum_{k\in [r]}p_{i,k}p_{i,k} -1,$$
which is the axiom of $\bbI(r,n)$ for $i=j$ after the substitution of the  variables.  The proof has degree $2$. 
The result follows immediately  from Lemma \ref{lem:PCred} and Theorem \ref{thm:PCdeglb}. 

\end{proof}

\begin{corollary}
Any $PC$ proof of $AB=I \vdash BA=I$, where $A,B$ are square $n\times n$ $\{0,1\}$ matrices requires degree $n/2+1$.
\end{corollary}

\begin{proof}
Follows immediately from Theorem~\ref{thm:rank} and Lemma~\ref{lem:inv_to_rank}.
\end{proof}
